% Discussion (DECISIONS §F37), open problems (§A10, §D1: every question is cited), acknowledgements and code and data
% (§A4, exact texts for Paper 2).
\section{Discussion}\label{sec:discussion}

\paragraph{What counts as one hypothesis.} Whether a steeper simplicity prior can separate a model from its noise depends on what counts as one hypothesis. When a universal mixture is a member of the class, as in SI and in SFI with shared randomness, it cannot for any model longer than $\ell^*_\lambda$, a bound that does not depend on the data, and for fixed $\lambda$ a computable teacher's labels, errors included, are eventually predicted (\cref{thm:absorb:absorption}). With private randomness at each input no member is universal, and steeper selection can work; as in algorithmic statistics (\cref{prop:algstat:parity}), the product view counts as model what is reused across inputs and as noise what is private to each input. This is a modelling choice with a price on shared latent causes (\cref{rem:absorb:why,prop:absorb:price}), and the protection ends when inputs carry labelled history (\cref{prop:absorb:history}).

\paragraph{The window and its price.} For error rules incompressible relative to the model, near-optimal product hypotheses are, up to explicit slack, the simple model with noise for $2^{-m}<c<c_{\rm hi}$ (\cref{thm:teacher:main}). The lower end is sharp because a bit of information about the error rule costs $c$ per sample in length and saves about $2^{-m}$ per sample in risk. So $\lambda$ must grow linearly with $n$, and the learner is then deliberately inconsistent for the teacher's labels: it discards every component whose rate is below $c$, valid or not, and its guaranteed code length exceeds a program's own loss by $c\,\ell(\pi)$ bits per sample (\cref{rem:limits:schedule}). Imitation data cannot justify this price, because strictly proper scores prefer the teacher's channel, errors included (\cref{prop:limits:noident}).

\paragraph{The gap to realistic errors.} Human labelling errors are often cheap, frequent patterns rather than incompressible sets. In the parity instance steeper selection provably fails for such errors in two ways: a fixed cheap rule defeats every $c$ once the simple model is long enough (\cref{prop:limits:rateblind,lem:limits:robust}), and a rule made of a cheap region plus random points, with the original density and description length, makes the good model non-optimal on about a third of the window (\cref{prop:limits:envelope}). Whether any window survives for structured errors in general is open. Steeper simplicity sorts components by rate, not by validity (\cref{prop:limits:separable}), so by itself it does not tell errors from rare valid rules. We do not expect steeper penalisation alone to solve the general problem of learning a norm from imperfect imitation; it works in one toy case and fails in others.

\paragraph{The model--randomness split.} At $\lambda=1$ a Bayes mixture is indifferent, up to factors that grow at most triply logarithmically in the amount of data, between a sampler and the programs that fill in its sampled bits (\cref{thm:split:lambda1}). For fixed $\lambda>1$ a steeper prior over product hypotheses breaks the tie in favour of samplers, in the limited sense of \cref{rem:split:reading}; in SI's representation by deterministic programs no form of steepness does (\cref{rem:split:siprograms}).

\paragraph{Relation to other work.} The $\lambda$-MAP is a two-part code with complexity weight $\lambda$, as in $\lambda$-weighted MDL in noisy classification, and maximises the tempered posterior, a generalised posterior with learning rate $1/\lambda$; the power-prior posterior has the same maximiser but is Bayes' rule with the prior $2^{-\lambda\ell}$. Those works aim at the data distribution; our target differs from it by design (\cref{sec:related}). A neural network whose labels are independent given its parameters has a Bayes posterior that mixes product laws, so \cref{thm:split:private} applies to countable discretisations; whether a parameter prior behaves like a description-length prior, as \cref{thm:teacher:main} would need, is the bridge question of the companion paper \citep[Question~10.9]{claudeSFI2026fundamentals}, which the slides of \citet{vaintrobhanni2026nnbayes} motivate (\cref{rem:related:nn}).

\subsection{Open problems}\label{sec:discussion:open}

\begin{enumerate}[label=(\arabic*),itemsep=1pt]
\item \emph{Structured error rules.} For envelopes whose cheap region has error density below one, only part (i) of \cref{prop:limits:envelope} carries over; whether near-optimal hypotheses memorise the random part, and what is selected above $c_{\rm env}$, are open (\cref{q:limits:open}).
\item \emph{Other model classes and flip rates.} Which simple-model classes give a window? Parities do; we expect that majority-like functions do not give a window. Is there a window for flip rates up to $1/4$, where the kink exists (\cref{q:teacher:open,q:limits:open})?
\item \emph{Non-uniform inputs and finite samples.} A version of \cref{thm:teacher:main} for non-uniform input laws (only the separable case is proved, \cref{prop:limits:inputlaw}); near-population constants at realistic sample sizes: they need about $10^{11}$ samples, and coarser guarantees hold from about $10^9$ (\cref{cor:teacher:sampled,q:teacher:open}).
\item \emph{Choosing $c$.} A principled choice of $c$ that does not use the imitation data, and the extra information that would tell errors from rare valid rules (\cref{q:limits:open}).
\item \emph{Shared hypotheses.} Is every lower-semicomputable law of a random partial function dominated by $\MF$ up to a constant (\cref{q:absorb:shared})?
\item \emph{Resource bounds and computability.} Does a time-bounded shared class contain a short member that absorbs everything at $\lambda>1$? Is there a computable map from data sets to $\lambda$-MAPs (\cref{q:related:open})?
\item \emph{Lower-semicomputable product models.} Does \cref{prop:algstat:coarse} hold for them (\cref{q:algstat:lsc})?
\end{enumerate}

\section*{Acknowledgements}
The coin example discussed in \cref{sec:split:lambda1} and in the companion paper is due to Sam Eisenstat (personal communication). The question of \cref{rem:related:nn} comes from joint work with Dmitry Vaintrob \citep{vaintrobhanni2026nnbayes}.

This paper was developed and written with extensive assistance from Claude (Anthropic), an AI system. Under the author's direction, Claude produced the proofs, the computations and the text; further, independent instances of Claude checked the proofs in two rounds and re-implemented the main computations, and the results were revised in response. The earlier report \citep{claude2026whatfollows}, whose results are re-derived here, was prepared in the same way. No human has yet independently checked all of the mathematics. The author takes responsibility for the content of this paper.

\section*{Code and data}
Every numerical statement comes from a deterministic or seeded Python script (Python~3.11 with NumPy and SciPy) whose output is saved next to it. The scripts, their outputs and independent re-implementations of the main computations are in the folder \texttt{solomonoff-function-induction/} of the repository \url{https://github.com/kaarelh/AI-works}; its README lists the script behind each number.
